\documentclass[10pt,a4paper]{amsart}
\usepackage[margin=19mm]{geometry}
\usepackage{amsmath,amssymb,mathtools}
\usepackage[scr=rsfs,cal=euler]{mathalfa}
\usepackage{xfrac}
\usepackage[unicode,hidelinks,bookmarks=false]{hyperref}
\newcommand{\ie}{i.e.\ }
\newcommand{\cf}{cf.\ }
\newcommand{\resp}{respectively\ }
\newcommand{\Subsection}{\subsection}
\providecommand{\nobreakdash}{\nobreak\hbox{-}}
\setlength{\emergencystretch}{2em}
\multlinegap0pt
\usepackage{mathtools}
\usepackage{mathrsfs}
\usepackage{appendix}
\newcommand{\joao}[1]{{\color{blue} #1}}
\newcommand{\discuss}[1]{{\color{orange} \textbf{[CHECK/DOUBT: #1]}}}
\newtheorem{theorem}{Theorem}[section]
\newtheorem{lemma}[theorem]{Lemma}
\newtheorem{corollary}[theorem]{Corollary}
\newtheorem{proposition}[theorem]{Proposition}
\newtheorem{definition}[theorem]{Definition}
\newtheorem{conjecture}[theorem]{Conjecture}
\newtheorem{remark}[theorem]{Remark}
\newtheorem{notation}[theorem]{Notation}
\newtheorem{example}[theorem]{Example}
\newtheorem*{problem}{Open problem}
\newcommand{\Weyl}{A_n(\mathbb{K})} % Álgebra de Weyl
\newcommand{\Invar}{A_n^G(\mathbb{K})}           % Subálgebra de Invariantes
\newcommand{\GKdim}{\text{GKdim}}    % Dimensão de Gelfand-Kirillov
\def\k{\mathbb K}
\def\C{\mathbb C}
\newcommand{\eq}{equation}
\newcommand{\Mat}{{\tt{Mat}}}
\newcommand{\Alg}{{\tt{Alg}}}
\newcommand{\HH}{{\rm{HH}}}
\newcommand{\Res}{{\rm{Res}}}
\newcommand{\Ann}{{\rm{Ann}}}
\newcommand{\Frac}{{\rm{Frac}}}
\newcommand{\Fract}{{\rm{Fract}}}
\newcommand{\Quot}{{\rm{Quot}}}
\newcommand{\eps}{\varepsilon}
\newcommand{\Tor}{{\rm{Tor}}}
\newcommand{\pr}{{\tt can\,}\,}
\newcommand{\ol} {\mbox{\rm{$1$-length}}}
\newcommand{\len} {\mbox{\rm{length}}}
\newcommand{\Spec}{{\rm{Spec}}}
\newcommand{\Span}{{\rm{span}}}
\newcommand{\Pic}{{\rm{Pic}_{\C}}}
\newcommand{\Aut}{{\rm{Aut}_{\C}}}
\newcommand{\Autk}{{\rm{Aut}}}
\newcommand{\Autg}{{\rm{Aut}_{\Gamma}}}
\newcommand{\id}{{\rm{Id}}}
\newcommand{\Der}{{\rm{Der}}}
\newcommand{\Div}{{\rm{Div}}}
\newcommand{\rk}{{\rm{rk}}}
\newcommand{\Tr}{{\rm{Tr}}}
\newcommand{\tr}{{\rm{tr}}}
\newcommand{\Ker}{{\rm{ker}}}
\newcommand{\Inn}{{\rm{Inn}}}
\newcommand{\diag}{{\rm{diag}}}
\newcommand{\Coker}{{\rm{Coker}}}
\newcommand{\ev}{{\rm{ev}}}
\newcommand{\im}{{\rm{im}}}
\newcommand{\balpha}{{\boldsymbol{\alpha}}}
\newcommand{\bmu}{{\boldsymbol{\mu}}}
\newcommand{\bn}{{\boldsymbol{n}}}
\newcommand{\bk}{{\boldsymbol{k}}}
\newcommand{\bU}{{\boldsymbol{U}}}
\newcommand{\bL}{{\boldsymbol{L}}}
\newcommand{\bM}{{\boldsymbol{M}}}
\newcommand{\be}{{\boldsymbol{\rm e}}}
\newcommand{\Ui}{U_{\infty}}
\newcommand{\blambda}{{\boldsymbol{\lambda}}}
\newcommand{\bgg}{{\boldsymbol{g}}}
\newcommand{\Vi}{V_{\infty}}
\newcommand{\vi}{v_{\infty}}
\newcommand{\ei}{e_{\infty}}
\newcommand{\lambdai}{\lambda_{\infty}}
\newcommand{\supp}{{\rm{supp}}}
\newcommand{\into}{\,\,\hookrightarrow\,\,}
\newcommand{\too}{\,\,\longrightarrow\,\,}
\newcommand{\onto}{\,\,\twoheadrightarrow\,\,}
\newcommand{\Sp}{{\rm{Sp}}}
\newcommand{\GL}{{\rm{GL}}}
\def\k{\mathbb K}
\def\I{\mathcal I}
\newcommand{\Tdeg}{{\rm Tdeg}}
\newcommand{\LD}{{\rm LD}}
\newcommand{\Hom}{{\rm Hom}}
\newcommand{\gr}{{\rm gr}}
\newcommand{\Supp}{{\rm Supp}}
\newcommand{\UGK}{{\rm UGK}}
\def\K{\mathcal K}
\def\L{\mathcal L}
\def\LL{\mathbf L}
\def\M{\mathcal M}
\def\Q{\bar{\mathbb Q}}
\def\D{\mathcal D}
\def\Dc{\mathcal{D}_c}
\def\Oo{\mathcal O}
\def\Z{\mathbb Z}
\def\R{\mathcal R}
\def\N{\mathbb N}
\def\g{\mathfrak g}
\def\h{\mathfrak h}
\def\tt{\mathfrak t}
\def\ad{\operatorname{ad}}
\def\aa{\mathfrak a}
\def\bb{\mathfrak b}
\def\AFC{\mathsf {AFC\_0}}
\def\AFCp{\mathsf {AFC\_p}}
\def\KK{\mathbb K}
\def\AA{\mathbb A}
\def\A{\mathcal{A}}
\def\AAA{\mathfrak A}
\def\XX{\mathfrak XX}
\def\G{\mathcal G}
\def\c{\mathfrak c}
\newcommand{\cfs}{\operatorname{cfs}}
\def\M{\mathcal M}
\def\Gskew{\mathcal{F}}
\def\Ginv{\mathcal{F}^*}
\def\H{\mathcal H}
\newcommand{\AMod}{A\text{-}\mathsf{Mod}}
\newcommand{\Amod}{A\text{-}\mathsf{mod}}
\newcommand{\Mod}{\mathsf{Mod}}
\newcommand{\HAHC}{\mathsf{Hopf}\text{-}\mathcal{AHC}}
\def\mod{\mathsf{mod}}
\def\supp{\mathsf{supp}}
\newcommand{\End}{\operatorname{End}}
\newcommand{\Za}{\mathbb{Z}^{\oplus \alpha}}
\newcommand{\SpecMax}{\operatorname{SpecMax}}
\begin{document}
\title[Irreducible non-holonomic modules for rational Cherednik alg]{Irreducible non-holonomic modules for rational Cherednik algebras}
\author{Felipe Albino dos Santos; João Schwarz}
\date{}
\maketitle
\noindent\emph{Source and licence.} Irreducible non-holonomic modules for rational Cherednik algebras. Version arXiv:2607.27038v1 (CC BY 4.0). This material is distributed under the Creative Commons Attribution 4.0 International licence, http://creativecommons.org/licenses/by/4.0/, as recorded in the arXiv OAI licence metadata for this exact version and shown on the abstract page https://arxiv.org/abs/2607.27038v1. Full rights notice: licenses/arxiv-2607.27038-CC-BY-4.0.txt.\par\smallskip
\noindent\textit{Editorial scope.} Counted unit: the original \texttt{theorem} environment \texttt{essential-invariant} at source lines 276--284 together with its complete proof at lines 285--287; the delivered document keeps source lines 196--294 so that every referenced label and every citation resolves inside it. Native editable source is the paper's own concatenated TeX; the AMS wrapper is editorial formatting only. No publisher class, upstream script or unrelated artwork is redistributed.
\section{Preliminaries}

\subsection{Gelfand-Kirillov dimension}
The Gelfand-Kirillov dimension is an important dimension in ring theory and is the basis of our ring-theoretic approach to the notion of holonomic modules, as discussed in the Introduction. Canonical references are \cite{KL}, \cite{Lorenz0}, and \cite[Chapter 8]{McConnell}.
\begin{definition}
Let $\Phi$ be the set of functions $f: \mathbb{N} \rightarrow \mathbb{R}$ such that there exists an $n_0 \in \mathbb{N}$ with:
\[f(n)>0\text{ and } f(n+1) \geq f(n), \quad\forall \, n \geq n_0.\]
For two functions in $\Phi$, we write $f \leq^* g$ if there are constants $c, m \in \mathbb{N}$ such that $f(n) \leq cg(mn)$. If $g \leq^* f$ as well, we write $f \sim g$. We represent the class of $f$ in $\Phi/\sim$ by $\mathcal{G}(f)$. This equivalence class of $f$ is called its \emph{growth}, and the relation induced by $ \leq^* $ in $\Phi/\sim$ is denoted $\leq$.
\end{definition}
\begin{notation}
Let $y \in \mathbb{R}, y \geq 0$. The growth of the function $n \mapsto n^y$ is denoted $\mathcal{P}_y$.
Let $y \in \mathbb{R}, y>0$. The growth of the function $n \mapsto e^{n^y}$ is $\mathcal{E}_y$.
\end{notation}
If $f(n)=\ln(n+1)$, then $\mathcal{G}(f) > \mathcal{P}_0$ but $\mathcal{G}(f) < \mathcal{P}_\varepsilon$ for any $\varepsilon >0$. So there is no analogue of the Archimedean property for this order.
Finally, we remark that not every pair $\mathcal{G}(f)$, $\mathcal{G}(h)$ is comparable.
\begin{notation}\label{asymptotic}
For $f \in \Phi$,
\[\gamma(f)= \limsup_{n \to \infty}\big(\log_n \, f(n)\big).\]
\end{notation}
\begin{proposition}\label{growth}
    Let $f,g \in \Phi$.
    \begin{enumerate}
        \item If $\mathcal{G}(f)=\mathcal{G}(g)$, then $\gamma(f)=\gamma(g)$.
        \item $\gamma(f+g)=\sup \{ \gamma(f), \gamma(g) \}$.
        \item If $f(n)=p(n)$ for $n\gg 0$ and $p(x) \in \mathbb{R}[x]$, then $\gamma(f)=\deg \, p$.
        \item If $\gamma(f)=\lim_{n \to \infty} \log_n \, f(n)$, then $\gamma(fg)= \gamma(f)+\gamma(g)$.
        \item If $g(n) \leq f(an+b), \, n\gg 0$, where $a,b \in \mathbb{N}$, then $\gamma(g) \leq \gamma(f)$.
    \end{enumerate}
\end{proposition}
\begin{proof}
    \cite[Lemma 2.1b)]{KL}, \cite[Lemma 8.1.7]{McConnell}.
\end{proof}
\begin{definition}
Let $A$ be an affine algebra and $V$ a finite-dimensional generating space for $A$. Let $d_V(n)=\dim \, \sum_{i=0}^n V^i$. Then $\mathcal{G}(d_V)$ is independent of the choice of the generating space $V$ (\cite[Lemma 1.1]{KL}). Hence,
\[ \gamma(d_V)=\limsup_{n \rightarrow \infty} \big(\log_n(d_V(n))\big),\]
is also independent of $V$, so by Proposition \ref{growth}(1), this number is well-defined and is called the \emph{Gelfand-Kirillov dimension} of $A$, denoted $\GKdim \, A$. We also put $\mathcal{G}(A)=\mathcal{G}(d_V)$, where $V$ is any choice of finite-dimensional generating space for $A$, and call it \emph{the growth of} $A$.
\end{definition}
\begin{remark}
When $1 \in V$, $d_V(n)$ reduces to $\dim \, V^n$, and $V$ is called a frame.
\end{remark}
\begin{example}
The Gelfand-Kirillov dimension of the Weyl algebra $\Weyl$ over a field $\k$ of zero characteristic is $2n$ \cite[Proposition 8.1.15(ii)]{McConnell}.
\end{example}
\begin{definition}
    Let $M$ be a finitely generated left $A$-module, where $A$ is an affine algebra. Choose a finite-dimensional space $F \subset M$ that generates it as a module, and let $V$ be a frame for $A$. Let $d_{F,V}(n)=\dim \, V^n F$. The \emph{Gelfand-Kirillov dimension} of $M$ is
    \[ \GKdim \, M=\GKdim_A\, M=\limsup_{n \to \infty} \, \log_n \, d_{F,V}(n), \]
    and the value is independent of the choices of $F, V$ by \cite[Proof of Lemma 1.1]{KL}.
\end{definition}
\begin{example}
Let $M:=\k[x_1,\ldots,x_n]$ be the polynomial module for $\Weyl$. Then its Gelfand-Kirillov dimension is $n$ (e.g., \cite[Theorem 8.5.10]{McConnell}).
\end{example}
We recall here two results about the Gelfand-Kirillov dimension that will be important for us.
\begin{proposition}\label{basic-GK-1}
    Let $\theta: S \rightarrow R$ be a homomorphism of algebras, which allows us to regard any left $R$-module $M$ as a left $S$-module. Then $\GKdim_S \, M \leq \GKdim_R \,  M$, with equality $\GKdim_S \, M = \GKdim_R \,  M$ when $R$ is a finitely generated $S$-module via $\theta$. \end{proposition}
\begin{proof}
    \cite[Proposition 1.6(d)]{Lorenz0}.
\end{proof}
In general, the Gelfand-Kirillov dimension behaves poorly with respect to localization \cite[Chapter 4]{KL}. For instance, while $\GKdim \, \Weyl=2n$, the Gelfand-Kirillov dimension of the Weyl field $\Frac \, \Weyl$ is $\infty$ (e.g., \cite[Theorem 8.17]{KL}). The following is one of the few known cases in which localization behaves well.
\begin{theorem}\label{basic-GK-2}
    Let $A$ be an affine $\k$-algebra and $S$ a multiplicatively closed subset of $A$ consisting of regular elements. Assume $S$ is commutative and, for each $s \in S$, $\operatorname{ad}(s)$ is a locally nilpotent derivation of $A$. Assume as well that $A$ has no $S$-torsion. Then $S$ is an Ore set in $A$, and for each $A$-module $M$, $\GKdim_A M= \GKdim_{AS^{-1}} MS^{-1}$.
\end{theorem}
\begin{proof}
    \cite[Theorem 3.2]{Lorenz0}.
\end{proof}
\subsection{Noncommutative invariant theory}
Let $G$ be a finite subgroup of $\Autk_\KK \, R$, where $R$ is a $\KK$-algebra. We have the \emph{trace function} $\tau:R\rightarrow R^G$, where $\tau(r)=\sum_{g \in G} r^g, \, r \in R$.
We begin by recording a simple lemma.
\begin{lemma}\label{most-basic-lemma}
Let $R$ be a $\KK$-algebra and $G \subset \Autk_\KK R$ a finite group such that $|G|^{-1} \in \KK$. Let $e=1 /|G| \sum_{g \in G} g$. Denoting the skew group ring as $R*G$, we have that $e^2=e$, $e(R*G)=eR$, $e(R*G)e=eR^G\simeq R^G$ and $ere=e\tau(r/|G|)$ for $r \in R$. More generally, if $\Lambda \subset R$ is a subalgebra with $G(\Lambda) \subset \Lambda$, then $\tau(\Lambda)=\Lambda^G$ and $e \Lambda e=e \Lambda^G \simeq \Lambda^G$.
\end{lemma}
\begin{proof}
    For the first half, \cite[Lemma 2.1 and its proof]{Montgomery}, with isomorphism $\psi: R^G \simeq eR^G$ given by $ r \mapsto er$, $ r \in R^G$, as $eR^G=R^Ge$. The second half follows from the first: $e \Lambda e= e \tau(\Lambda)= e \Lambda^G$, by hypothesis, and the restriction of $\psi$ to $ \Lambda^G$ gives us the desired result.
\end{proof}
\begin{theorem}\label{MS}
Let $R$ be a finitely generated Noetherian $\k$-algebra. If $G$ is a finite group of automorphisms of $R$, with $|G|^{-1} \in \k$, then $R^G$ is a finitely generated $\k$-algebra, $R^G$ is Noetherian, and $R$ is a finitely generated $R^G$-module.
\end{theorem}
\begin{proof}
    The first claim is proven in \cite{MS}, the second one in \cite[Corollary 1.12]{Montgomery}, and the third one in \cite[Corollary 5.9]{Montgomery}.
\end{proof}
When our algebra is simple, more can be said:

\begin{theorem}\label{essential-invariant}
    Let $R$ be a simple $\k$-algebra, and $G$ a finite group of outer automorphisms of $R$ such that $|G|^{-1} \in \k$. Then
    \begin{enumerate}
        \item $R^G$ is a simple ring.
        \item $R$ is a finitely generated projective $R^G$-module.
        \item $R^G$ and $R*G$ are Morita equivalent. In particular, if $R$ is Noetherian, $R^G$ is Noetherian.
%        \item $t(R*g)$
    \end{enumerate}
\end{theorem}

\begin{proof}
    \cite[Theorem 2.4, Theorem 2.5, Corollary 2.6]{Montgomery}.
\end{proof}

This is relevant to the study of $\Weyl$:
\begin{proposition}\label{Weyl-outer}
    Every automorphism of the Weyl algebra is outer.
\end{proposition}
\begin{proof}
    \cite[2.4.1]{Dumas}.
\end{proof}

\begin{thebibliography}{99}
\bibitem[Dum06]{Dumas} {\sc F. Dumas}, {\em An introduction to noncommutative polynomial invariants}. Lecture notes from the conference ``Homological methods and representations of non-commutative algebras'', Mar del Plata, Argentina March 6--17, 2006.\\ \href{https://lmbp.uca.fr/$\sim$fdumas/recherche.html}{https://lmbp.uca.fr/$\sim$fdumas/recherche.html}
\bibitem[KL00]{KL} {\sc G. R. Krause and T. H. Lenegan}, {\em Growth of Algebras and Gelfand-Kirillov Dimension}, Grad. Stud. Math., vol. 22, American Mathematical Society, Providence, RI, 2000.%\\ %\href{https://doi.org/10.1090/gsm/022}{https://doi.org/10.1090/gsm/022}
\bibitem[Lor88]{Lorenz0} {\sc M. Lorenz}, {\em Gelfand-Kirillov dimension and Poincaré series}, Cuadernos de Algebra, vol. 7, Universidad de Granada, Granada, España, 1988. %\href{https://www.math.temple.edu/~lorenz/pubs.html}{https://www.math.temple.edu/$\sim$lorenz/pubs.html}
\bibitem[MS81]{MS} {\sc S. Montgomery, L. W. Small}, {\em Fixed rings of noetherian rings}, Bull. London. Math. Soc. {\bf 13}(1) (1981) 33--38.%\\ %\href{https://doi.org/10.1112/blms/13.1.33}{https://doi.org/10.1112/blms/13.1.33}
\bibitem[MR01]{McConnell} {\sc J. C. McConnell and J. C. Robson} {\em Noncommutative Noetherian rings}, Grad. Stud. Math., vol. 30, American Mathematical Society, Providence, RI, 2001.%\\ %\href{https://doi.org/10.1090/gsm/030}{https://doi.org/10.1090/gsm/030}
\bibitem[Mon80]{Montgomery} {\sc S. Montgomery}, {\em Fixed rings of finite automorphism groups of associative rings}, Lecture Notes in Math., vol. 818, Springer, Berlin, 1980.%\\ %\href{http://doi.org/10.1007/BFb0091561}{http://doi.org/10.1007/BFb0091561}
\end{thebibliography}
\end{document}
